\section*{Devoir maison : Diviser pour Régner (Le Tri Fusion)}

Maki Tire-Vite, le cow-boy du lagon, doit trier une grande file de Lagonnet. Pour aller plus vite, il utilise une méthode redoutable : le \textbf{tri fusion}. Il sépare la file en deux, demande à ses adjoints de trier chaque moitié, puis il \textbf{fusionne} les deux moitiés triées pour n'en former qu'une seule, parfaite.

\begin{lstlisting}[language=Python]
def fusion(gauche, droite):
    resultat = []
    i = j = 0
    # On compare les premiers de chaque file et on prend le plus petit
    while i < len(gauche) and j < len(droite):
        if gauche[i] < droite[j]:
            resultat.append(gauche[i])
            i += 1
        else:
            resultat.append(droite[j])
            j += 1
    # On ajoute les restes des files
    resultat.extend(gauche[i:])
    resultat.extend(droite[j:])
    return resultat

def tri_fusion(tab):
    if len(tab) <= 1:
        return tab
    milieu = len(tab) // 2
    gauche = tri_fusion(tab[:milieu])
    droite = tri_fusion(tab[milieu:])
    return fusion(gauche, droite)
\end{lstlisting}

\begin{Exercise}[title={La stratégie de Maki Tire-Vite, le cow-boy du lagon}]

\Question \textbf{L'étape de Fusion.} 
On suppose que Maki Tire-Vite, le cow-boy du lagon, a déjà deux groupes triés : \texttt{gauche = [0.9, 1.5]} (Petit-Îlot et Grande-Barge) et \texttt{droite = [1.2, 2.1]} (Moyen-Maki et Mont-Choungui).
Complétez les étapes de la fonction \texttt{fusion} :
\begin{itemize}
    \item On compare 0.9 et 1.2. Le plus petit est ........ On l'ajoute au résultat.
    \item On compare 1.5 et 1.2. Le plus petit est ........ On l'ajoute au résultat.
    \item On compare 1.5 et 2.1. Le plus petit est ........ On l'ajoute au résultat.
    \item Il ne reste que 2.1, on l'ajoute à la fin.
    \item Résultat final : \texttt{[ \quad , \quad , \quad , \quad ]}
\end{itemize}

\Question \textbf{Simulation visuelle.}
Dessinez l'arbre de division pour le tableau \texttt{tab = [1.5, 0.9, 2.1, 1.2]} en suivant l'exemple :



\begin{center}
\begin{tabular}{cc}
      & [1.5, 0.9, 2.1, 1.2] \\
     / & \textbackslash \\
    {[1.5, 0.9]} & {[2.1, 1.2]} \\
   / \quad \textbackslash & / \quad \textbackslash \\
  {[1.5]} \quad [0.9] & {[2.1]} \quad [1.2] \\
   \hline
   \multicolumn{2}{c}{\textit{(Maintenant, on remonte en fusionnant...)}} \\
   \textbackslash \quad / & \textbackslash \quad / \\
   {[0.9, 1.5]} & {[..........]} \\
    \textbackslash & / \\
    \multicolumn{2}{c}{[0.9, 1.2, 1.5, 2.1]} \\
\end{tabular}
\end{center}

\Question \textbf{Analyse de la complexité.}
Contrairement au tri par insertion, le tri fusion coupe toujours le tableau en deux, peu importe si les Lagonnet sont déjà rangés ou non.
\begin{itemize}
    \item Est-ce que le temps de calcul change beaucoup si le tableau est déjà trié ? ....................
    \item Si on a 1000 Lagonnet, le tri fusion est beaucoup plus rapide que le tri par insertion. On dit que sa complexité est en $n \log(n)$. À votre avis, pourquoi "couper en deux" aide-t-il l'ordinateur ? \\
    ..................................................................................................................
\end{itemize}

\Question \textbf{Le cas de base.}
Dans la fonction \texttt{tri\_fusion}, à quoi sert la ligne \texttt{if len(tab) <= 1: return tab} ? \\
\textit{(Indice : peut-on couper en deux un Lagonnet tout seul ?)} \\
..................................................................................................................

\end{Exercise}